\begin{abstract}
\noindent This experiment focuses on implementing line-tracking behavior for a robotic car, with the primary objective of enabling it to perform precise path-following tasks based on infrared (IR) sensor feedback.
Specifically, the study investigates the variation in analog-to-digital converter (ADC) values of incoming IR signals as the car moves at different distances from the guiding line.
By analyzing these signal changes, we aim to establish a quantitative relationship between the IR sensor outputs and the car’s lateral position relative to the path.
Furthermore, the experiment explores how different line shapes affect the IR sensor’s performance, thereby providing insight into how environmental factors influence the accuracy and stability of line tracking.
The results offer foundational data to support autonomous navigation of the car.

\end{abstract}